%% ========================================================================
%% Bhaskara II: Bijaganita (Part 3: pages 115-170)
%% BILINGUAL English-Sanskrit Edition
%% Parallel columns: Sanskrit (Devanagari) | English translation with IAST
%% ========================================================================
\documentclass[11pt,a4paper]{article}

%% -------- Core Packages --------
\usepackage{fontspec}
\usepackage{polyglossia}
\usepackage{amsmath,amssymb,amsthm}
\usepackage{geometry}
\usepackage{graphicx}
\usepackage{xcolor}
\usepackage{titlesec}
\usepackage{fancyhdr}
\usepackage{enumitem}
\usepackage{array,booktabs,longtable,multirow}
\usepackage{hyperref}
% verse, lettrine, microtype, tocloft removed for compilation
\usepackage{indentfirst}

%% -------- Page Geometry --------
\geometry{
  paperwidth=210mm,paperheight=297mm,
  left=18mm,right=18mm,top=28mm,bottom=28mm,
  headheight=14pt,footskip=30pt
}

%% -------- Language Setup --------
\setmainlanguage{english}
\setotherlanguage{sanskrit}

%% -------- Font Configuration --------
\setmainfont{Linux Libertine O}
\setsansfont{Linux Biolinum O}
\setmonofont{DejaVu Sans Mono}

%% Devanagari (Sanskrit/Hindi) Font
\newfontfamily\devanagarifont[Script=Devanagari,Path=/usr/share/fonts/truetype/noto/]{NotoSerifDevanagari-Regular.ttf}
\newfontfamily\devanagarifontsf[Script=Devanagari,Path=/usr/share/fonts/truetype/noto/]{NotoSansDevanagari-Regular.ttf}

%% -------- Colors --------
\definecolor{titlecolor}{RGB}{45,55,72}
\definecolor{sectioncolor}{RGB}{55,65,81}
\definecolor{notecolor}{RGB}{120,53,15}
\definecolor{rulecolor}{RGB}{180,160,140}
\definecolor{sanskrcolor}{RGB}{60,30,10}
\definecolor{engcolor}{RGB}{30,50,80}

%% -------- Section Formatting --------
\titleformat{\section}
  {\normalfont\Large\bfseries\color{sectioncolor}}
  {\thesection}{1em}{}
\titleformat{\subsection}
  {\normalfont\large\bfseries\color{sectioncolor}}
  {\thesubsection}{1em}{}
\titleformat{\subsubsection}
  {\normalfont\normalsize\bfseries\color{sectioncolor}}
  {\thesubsubsection}{1em}{}

%% -------- Header/Footer --------
\pagestyle{fancy}
\fancyhf{}
\fancyhead[L]{\small\textit{Bijaganita -- Bilingual Edition (Sanskrit-English)}}
\fancyhead[R]{\small\thepage}
\renewcommand{\headrulewidth}{0.4pt}
\renewcommand{\footrulewidth}{0pt}

%% -------- Custom Commands --------
\newcommand{\longline}{\noindent\rule{\textwidth}{0.4pt}}
\newcommand{\shortline}{\noindent\rule{0.5\textwidth}{0.4pt}\par}
\newcommand{\cnote}[1]{\textcolor{notecolor}{\textit{#1}}}
\newcommand{\dev}[1]{{\devanagarifont #1}}

%% Bilingual parallel columns
\newcommand{\bicol}[2]{%
\noindent%
\begin{minipage}[t]{0.48\textwidth}
\begin{flushleft}
\small
#1
\end{flushleft}
\end{minipage}%
\hfill%
\begin{minipage}[t]{0.48\textwidth}
\begin{flushleft}
\small
#2
\end{flushleft}
\end{minipage}%
\vspace{0.4em}
}

%% Sanskrit verse environment (replacement for verse package)
\newenvironment{sanskritverse}{%
  \begin{quote}
}{%
  \end{quote}
}

%% -------- Hyperref Setup --------
\hypersetup{
  colorlinks=true,
  linkcolor=sectioncolor,
  citecolor=sectioncolor,
  urlcolor=sectioncolor,
  pdfencoding=auto
}

%% -------- TOC Depth --------
\setcounter{tocdepth}{3}
\setcounter{secnumdepth}{3}

%% ========================================================================
\newcommand{\vs}{\vspace{0.5em}}
\begin{document}

%% ========================================================================
%% Title Page
%% ========================================================================
\begin{titlepage}
\centering
\vspace*{1.5cm}

{\Huge\bfseries\color{titlecolor} Bījagaṇita\par}
\vspace{0.5cm}
{\Large\itshape Algebra of Bhāskara II\par}
\vspace{1.5cm}

{\LARGE\bfseries Part III: Pages 115--170\par}
\vspace{0.3cm}
{\large \textsc{Bilingual English--Sanskrit Edition}\par}
\vspace{0.5cm}
{\normalsize (Parts 7--9 of 9)\par}
\vspace{0.8cm}

\shortline
\vspace{0.5cm}

{\large\dev{श्रीभास्कराचार्यविरचितं बीजगणितम्}\par}
\vspace{0.3cm}
{\normalsize \textit{Śrībhāskarācāryaviracitaṃ Bījagaṇitam}\par}
\vspace{0.5cm}
{\normalsize The Bījagaṇita (Algebra) composed by Śrī Bhāskarācārya\par}

\vspace{0.8cm}
\shortline
\vspace{0.5cm}

{\normalsize
\begin{minipage}{0.85\textwidth}
\centering
\textbf{Parallel-Column Layout:}\\[0.3em]
\textsc{Left:} Sanskrit original in Devanagari script\\[0.2em]
\textsc{Right:} English translation with IAST transliteration of technical terms\\[0.5em]
Covering advanced topics in \textit{Varga-prakṛti} (Pell's equation),\\
\textit{Aneka-varṇa-madhyamāharaṇa-bhedāḥ} (equations with multiple unknowns),\\
\textit{Bhāvitam} (product equations), and concluding verses
\end{minipage}
\par}

\vfill

{\small\textit{Typeset from the scanned edition published by}}\\[0.3em]
{\small Chowkhamba Sanskrit Series, Varanasi\par}
\vspace{1cm}
{\small Modern scholarly typesetting\par}
{\footnotesize XeLaTeX with Devanagari support\par}

\end{titlepage}

%% ========================================================================
%% Table of Contents
%% ========================================================================
\tableofcontents
\newpage

%% ========================================================================
%% Introduction
%% ========================================================================
\section*{Introduction}
\addcontentsline{toc}{section}{Introduction}

This volume presents the final portion of Bhāskara II's \textit{Bījagaṇita}, comprising pages 115--170 of the original Chowkhamba Sanskrit Series edition. This section covers the most advanced topics in the treatise:

\begin{enumerate}[label=\arabic*.]
  \item \textbf{Aneka-varṇa-madhyamāharaṇa-bhedāḥ} (Equations with Multiple Unknowns) --- systematic methods for solving equations involving several unknown quantities, building on the single-variable equations treated earlier.
  
  \item \textbf{Bhāvitam} (Product Equations) --- equations involving products of unknowns, including the famous ``Bhāvita'' (product) class of problems.
  
  \item \textbf{Śeṣaṇa} and \textbf{Prakṣepaka-vidhi} --- special algebraic methods for remainder problems and the method of addition/substitution.
  
  \item \textbf{Concluding Verses} --- the epilogue (\textit{upasaṃhāra}) in which Bhāskara summarizes his work and pays homage to tradition.
\end{enumerate}

Bhāskara employs his characteristic notation where unknown quantities are denoted by the initial syllable \dev{याव} (\textit{yāva-}, from \textit{yāvattāvat}, ``as much as''), and coefficients are prefixed to these symbols. The final sections also include a publisher's catalogue listing the many Sanskrit texts available from the Chowkhamba Sanskrit Book Depot, Varanasi.

\longline

%% ========================================================================
%% Part 7: Pages 115-133
%% ========================================================================
\section{From the Varga-prakṛti Section (p.\ 115)}

The opening of this portion continues the discussion of the Varga-prakṛti (square-nature) equation, i.e., equations of the form $Nx^2 + k = y^2$. The text presents further rules and examples for finding solutions to such indeterminate equations.

\subsection*{Continuation of the Square-nature Method}

\begin{sanskritverse}
\dev{%
समशोधने हृते पक्षौ यावच्चावत्तावकस्य रूपं प्रक्षिप्य प्रथमपक्षमूलं या २ रू १ ।\\
परपक्षस्य काव १२ रू १ । वर्गप्रकृत्या मूले क २ उचे ७ वा क २८ उचे ९७ ।\\
कनिष्ठं कालकमानम् । उच्चिष्टस्य या २ रू १ समं हृत्वा लब्धं यावच्चावकमानम् ३ वा ४८ ।\\
स्वच्छानेनोत्थापने हृते जाती राशी १, ५ वा २०, ७६ इत्यादि ॥}
\end{sanskritverse}

\cnote{After equal subtraction, the two sides being reduced, adding the original constant to the first side, its root is $\sqrt{4y^2 + 1}$. For the second side: $\sqrt{12y^2 - 1}$. By the square-nature [method], the roots are $(2, 7)$ or $(28, 97)$. The lesser is the time-value. The residue from $\sqrt{4y^2 + 1}$ gives the yāvattāvat value 3 or 48. By self-substitution and reduction, the series of values 1, 5 or 20, 76, etc.}

\bigskip

\begin{sanskritverse}
\dev{%
अथापरं सूत्रं साध्येऽस्तु ॥\\[0.5em]
(१) द्वितीयपक्षे सति सममेव तु हत्यापवर्त्योन पदे प्रसाध्ये ।\\
उच्चिष्टं कनिष्ठेन तदा निहत्याच्छेदं वर्गेण हतोऽपवर्तः ॥ ६ ॥\\
कनिष्ठवर्गेण तदा निहत्योच्चिष्टं ततः पूर्ववदेव शेषम् ।\\
इत्यर्थात् ॥}
\end{sanskritverse}

\cnote{Rule 6: ``When the second side exists [and is not unity], having made [the two sides] equal, multiplying and reducing, the remainder [is obtained]; having multiplied the residue by the lesser [root], dividing by the square [nature], the reduced [equation is obtained]. Then having multiplied by the square of the lesser [root], the residue --- from that, as before, the remainder.'' Thus the meaning.}

\subsection*{Worked Example (p.\ 115)}

\begin{sanskritverse}
\dev{%
(१) वि०—कल्प्यते समी पक्षौ\\[0.3em]
\hspace*{2em} काव १\\
\hspace*{2em} यावव्। ई १ याव।ई १}
\end{sanskritverse}

\dev{%
अत्र प्रथमपक्षस्य मूले क १ द्वितीयपक्षस्य यावव्।ई १ याव।ई १ मूलेन सममिति ।\\
तत्र द्वितीयपक्षस्य मूले च का।}

$$\text{का} = \sqrt{\text{यावव।ई १ याव।ई १}} = \text{या}\sqrt{\text{याव।ई १ ई १}}$$

\dev{%
अत इदं याव।ई १ ई १ मूलदं ततो वर्गप्रकृतिविषयो यथा का वर्गः ई, गुणः ई, युतो मूलदं इति हत्वा यावच्चावकमानं उच्चिष्टं चास्य याव।ई १ ई १ मूलेन सममिति।}

\subsection*{Summary Rule (p.\ 115 bottom)}

The section concludes with a general procedure:

\begin{sanskritverse}
\dev{%
एवं बहुधा बुद्धिमत्त्वविचिन्त्यामिति सर्वमुपपन्नम् ॥}
\end{sanskritverse}

\cnote{``Thus in many ways, by the thoughtful application of intelligence, all [is demonstrated].'' Everything is proved.}

%% ========================================================================
%% Part 7 Section 2: Aneka-varna-madhyamaharana-bhedah
%% ========================================================================
\section{Aneka-varṇa-madhyamāharaṇa-bhedāḥ (pp.\ 116--133)}
\label{sec:aneka-varna}

This is the largest section of the final portion, dealing systematically with equations involving multiple unknown quantities (\textit{aneka-varṇa}, ``many colours''). The Sanskrit term \textit{varṇa} (colour) is the traditional designation for an unknown quantity in Indian algebra --- different unknowns being assigned different colours.

\subsection{General Rules for Multiple Unknowns (pp.\ 116--117)}

\begin{sanskritverse}
\dev{%
अथ सूत्रं साध्येऽस्तु ॥\\[0.5em]
(१) द्वितीयपक्षे सति सममेव तु हत्यापवर्त्योन पदे प्रसाध्ये ।\\
उच्चिष्टं कनिष्ठेन तदा निहत्याच्छेदं वर्गेण हतोऽपवर्तः ॥ ६ ॥\\[0.5em]
कनिष्ठवर्गेण तदा निहत्योच्चिष्टं ततः पूर्ववदेव शेषम् ।\\
इत्यर्थात् ॥}
\end{sanskritverse}

\subsection*{Worked Example 1 (p.\ 117)}

\begin{sanskritverse}
\dev{%
उदाहरणम् ॥\\[0.3em]
यस्य वर्गेऽद्वितिः पञ्चगुणा वर्गशतोनिता ।\\
मूलद्वा जायते राशिः गणितज्ञ वदाशु तम् ॥ १ ॥}
\end{sanskritverse}

\cnote{Example 1: ``A number whose square minus twice [itself], diminished by a hundred times five [i.e., 500], has two roots --- tell me that number quickly, O mathematician.''}

\dev{%
अत्र राशिः=या १ । अस्य वर्गेऽद्वितिः पञ्चगुणा वर्गशतोनिता यावव ५ याव १०० । अयं वर्ग इति कालकवर्गसमं हृत्वा गृहीतं कालकवर्गस्य मूलं का १ । द्वितीयपक्षस्य यावव ५ याव १०० ।\\
यावच्चावकगुणापवत्यं वर्गप्रकृत्या मूले क १० उचे २० वा क १७० उचे ३८० ।}

\subsection*{Worked Example 2 (p.\ 117)}

\begin{sanskritverse}
\dev{%
उदाहरणम् ॥\\[0.3em]
कयोः स्यादन्तरे वर्गौ वर्गयोगो ययोर्घनः ।\\
तौ राशी कथयामित्रो वहुधा बोज्ञवित्तम ॥ २ ॥}
\end{sanskritverse}

\cnote{Example 2: ``Two numbers whose difference of squares [equals] the sum of their squares --- tell me those numbers, O friend, in many ways, O best of algebraists.''}

\dev{%
अत्र राशी या १, का १ । अनयोरन्तरं या १ का १ नीलकवर्गसमं हृत्वा लब्धं यावच्चावकमानं का १ नीलं ई १ । अनेन यावच्चावकद्वितीयाप्य जाती राशी का १ नीलं ई, का १ ।\\
अनयोर्वर्गयोगः काव २ नीलका २ नीलव १ । एष घन इति नीलकवर्गयनसमं हृत्वा शोधने हृते जाते प्रथमपक्षे नीलव १ नीलव १ । द्वितीयपक्षे काव २ नीलका २ ।}

\subsection*{Rule for Three Unknowns (pp.\ 119--120)}

\begin{sanskritverse}
\dev{%
(१) साध्यकल्पो यदि वर्गवर्गेषद्वित्र्यादिवर्णस्य हते समं तम् ।\\
हृत्वा पदं तस्य तदन्वपक्षे वर्गप्रकृत्यैवकदैव मूले ॥\\
कनिष्ठमाच्छाद्य पदेन तुल्यं उच्चिष्टं द्वितीयेन समं विधायात् ॥ ८ ॥}
\end{sanskritverse}

\cnote{Rule 8: ``If the square of the square of an assumed [quantity], diminished by the square of one, two, three, etc., unknown quantities, having made [the two sides] equal, taking the square root of that, then by the square-nature method, the roots [are found]; covering the lesser root by an equal root, making the residue equal to the second [unknown]...''}

\subsection*{Worked Example with Three Unknowns (p.\ 121)}

\begin{sanskritverse}
\dev{%
उदाहरणम् ॥\\[0.3em]
त्रिकादिचतुरश्रेणां नच्छेदं क्वापि च यत् कलम् ।\\
तदेव त्रिगुणं क्षेत्रमष्टयुगच्छेदं भवेद्व्र ॥ ९ ॥}
\end{sanskritverse}

\dev{%
अत्र श्रेण्योन्या इः । आदि=३, चयः=२, गच्छः=या १ । आदि=३, चयः=२, गच्छः=का १ ।\\
अनयोः (क) कलं=याव १ या २, काव १ का २ । अनयोरन्यत् त्रिगुणं परस्मिं हृत्वा शोधनार्थं}

$$\text{न्यासः} \longrightarrow \text{याव ३ या ६}$$
$$\text{काव १ का २}$$

\dev{%
शोधने हृते पक्षौ त्रिगुणीहृत्व नव प्रक्षिप्य प्रथमपक्षस्य मूलं या ३ रू ३ ।\\
द्वितीयपक्षस्य काव ३ का ६ रू ९ । नीलकवर्गेण}

\subsection*{Rule for Equal-square Unknowns (p.\ 121)}

\begin{sanskritverse}
\dev{%
(९) प्रथमपक्षस्यद्विगुणीतच्छेदस्य मूलं नीलं प्रकल्प्य तद्वर्गेण समं परं पक्षं\\
हृत्वा पूर्वोक्तर्यस्य वासना चातिसरति ॥}
\end{sanskritverse}

\cnote{``Having assumed the root of twice the divisor of the first side as the `blue' [unknown], having divided the other side by its square, the desire is satisfied more than previously stated.''}

\subsection{The Method of Substitution (pp.\ 125--128)}

\begin{sanskritverse}
\dev{%
अथ सूत्रं सूचयम् ॥\\[0.5em]
(१) सकृपमयकमल्पकं वा विभोगमूलं प्रथमं प्रकल्प्य ।\\
योगान्तरशेषकमाजिताच्छेदान्तरशेषकतः पदं स्यात् ॥ १२ ॥}
\end{sanskritverse}

\cnote{Rule 12: ``Having assumed the root of the difference of [the two sides] as the first [unknown], or a small quantity, [and] the residue of the difference of sums as the square [of another unknown], the root [is obtained].''}

\dev{%
(९) वि०—अत्र कल्प्यते योगान्तरशेषमानम् = क्षे\\
वर्गान्तरशेषमानम् = क्षे$_2$, वर्गयोगशेषमानम् = क्षे$_3$\\
विभोगमूलम् = या, योगमूलम् = का}

$$\text{तदा प्रदर्शनार्थं विभोगः} = \text{या}^2 - \text{क्षे}_1, \text{ योगः} = \text{का}^2 - \text{क्षे}_1$$

$$\text{अल्पराशिः} = \frac{\text{का}^2 - \text{या}^2}{2}, \quad \text{वृहद्राशिः} = \frac{\text{का}^2 + \text{या}^2 - 2\text{क्षे}_1}{2}$$

\subsection*{Rule 13 (p.\ 129)}

\begin{sanskritverse}
\dev{%
शेषं ततः शेषकमूलमन्यच्च मूले विद्ध्वादसकृत् सम्मते ॥ ९ ॥\\
समाविते वर्गकृती तु यत्र तन्मूलमादाय च शेषकस्य ।\\
इष्टोद्धृतस्येर्विवर्जितस्य दलैन समं तद्विह तदैव कार्यम् ॥ १० ॥}
\end{sanskritverse}

\subsection*{Worked Example on the Product Method (p.\ 130)}

\begin{sanskritverse}
\dev{%
उदाहरणम् ॥\\[0.3em]
तौ राशी वद यत्कृत्योः समस्तगुणयोर्युतिः ।\\
मूलद्वा स्याद्‌योगमष्टौ मूलद्वा रूपसंयुतः ॥ १ ॥}
\end{sanskritverse}

\dev{%
अत्र राशी या १, का १ । अनयोर्वर्गयोः समस्तगुणयोर्युतिः यावव ७ काव ८ । अयं वर्ग इति नीलकवर्गसमं हृत्वा पक्षयोः कालकधनं प्रक्षिप्य नीलकपक्षस्य मूलं नी १ । परपक्षस्य याव ७ काव ८ ।}

\subsection*{Summary Verses (p.\ 131)}

\begin{sanskritverse}
\dev{%
अथ सूत्रं सूचयम् ॥\\[0.5em]
(१) सकृपमयकमल्पकं वा विभोगमूलं प्रथमं प्रकल्प्य ।\\
योगान्तरशेषकमाजिताच्छेदान्तरशेषकतः पदं स्यात् ॥ १२ ॥\\[0.5em]
वर्गान्तरशेषकसंमितिर्युता क्षेपेण हृत्वोर्ध्वितेन तै ततः ।\\
द्विधा पदं तत्पदयुगविभोगाज् मूलं युतमूलमतस्तयोरिति ॥}
\end{sanskritverse}

\cnote{``The difference of squares having a residue equal to the divisor, divided by the additive, then by the quotient, [the root] divided in two; from the difference of the pair of roots, the root [of the lesser] from the sum of the roots, thus [proceeding]...''}

%% ========================================================================
%% Part 8: Pages 134-152
%% ========================================================================
\section{Bhāvitam and Prakṣepaka-vidhi (pp.\ 134--152)}
\label{sec:bhavitam}

\subsection{Equations with Product Terms (pp.\ 134--138)}

This section deals with \textit{bhāvita} (product) equations, where the unknowns appear in product form. The method extends the techniques for multiple unknowns to handle more complex algebraic structures.

\begin{sanskritverse}
\dev{%
अथ वा स्वर्णकल्पनायां मन्दबुद्धीनां पुनः किञ्चित् पठितः ।\\
तत् सूचयामि ॥\\[0.5em]
हरमका यस्य हृति सूक्ष्मति सोऽपि द्वितीयपदयुतातः ।\\
तेनाहतोऽस्यवर्णो रूपपदेनान्वितः कल्पः ॥ १६ ॥}
\end{sanskritverse}

\cnote{Rule 16: ``The divisor multiplied by [some quantity] whose [result] is fine, added to the second term; having multiplied by that, the unknown [is obtained], combined with a root term [as] the assumption.''}

\subsection*{Worked Example (p.\ 137)}

\begin{sanskritverse}
\dev{%
उदाहरणम् ॥\\[0.3em]
न यदि पदं रूपाणां शिष्येत्तं तेन हरतोऽपि ।\\
नावच्छद्‌गुणो भवति न चेदेवमपि किल तथै ॥ १७ ॥}
\end{sanskritverse}

\dev{%
हरमतिः । यस्याङ्कस्य हृतिहरमका सती सूक्ष्यतीति निश्शेषा भवति अपि च सोऽप्यङ्को हर्यो रूपपदेन च गुणितो हरमकः सन् सूक्ष्यति तदा तेनाहतोऽस्यवर्णस्तेन रूपाणां घनमूलं न लभ्यते तदा तेन रूपेषु हरतोऽपि तावद्‌गुणो भवेत् तन्मूलं रूपपदं स्यात् ।}

\subsection{The Bhāvita Rule (pp.\ 140--141)}

\begin{sanskritverse}
\dev{%
भावितम् ॥\\[0.5em]
वर्णद्वयातिकूपेऽस्य मकरवेष्टनैकतः ।\\
एताम्या संयुताबुद्धि कल्पयो र्वैच्छया च तौ ॥ ३ ॥\\[0.5em]
वर्णद्वयो वर्गयोरन्ते क्रियते ते विपर्ययात् ।\\
समयोः पक्षयोरेकमाध्यमिकमपास्य तौ वर्णौ रूपाणि च}
\end{sanskritverse}

\cnote{The section on ``Product'': ``Two unknowns, their product having a coefficient, from one side having a single enclosure [i.e., bracketed], having assumed any sum of their values, [proceeding] with the pair of unknowns reversed, [taking] the square of two unknowns, the means [being] equal, having removed one middle term, the two unknowns and the roots...''}

\subsection*{Worked Example on Bhāvita (p.\ 142)}

\begin{sanskritverse}
\dev{%
उदाहरणम् ॥\\[0.3em]
चतुःषष्टिगुणयो राश्योः संयुतिर्द्वियुता तयोः ।\\
राशिघातेन तुल्येति ॥}
\end{sanskritverse}

\dev{%
तत्र यथोक्तं हृते पक्षौ $\left\{ \begin{array}{l} \text{या ४ का ३ रू २} \\ \text{या।का।भा १} \end{array} \right.$}

\dev{%
वर्णद्वयातिकूपेऽस्य १४ एतदेकेनाहतं हते जाते द्विके १, १४ । एते वर्णद्वया ४, ३ स्वैच्छया युते जाते यावच्चावककालकमाने ५, १८ वा १७, ५ । द्विकेन ५, ११ वा १०, ६ ॥}

\subsection{The Method of Mean Elimination (pp.\ 145--147)}

\begin{sanskritverse}
\dev{%
माध्यमाविगमः ॥\\[0.5em]
आसन्नमानस्य हराशमानम् अप्राप्तिगुण्ये सहिते क्षेपेण ।\\
पूर्णे स्याताम् अन्तरराशिकाभ्यां तदा हरांशो भवतीश्रमस्य ॥ १ ॥\\[0.5em]
आसन्नमानयोरन्तरमेव भवेत् ।\\
अंशस्थाने सदा रूपं चिन्त्येत तथा धीमता ॥ २ ॥\\
सर्वेषामासन्नमानेषु हरांशो भवति षट् ।\\
तथोत्तरोत्तरं सूक्ष्मतामासन्नानि भवन्ति हि ॥ ३ ॥}
\end{sanskritverse}

\cnote{The method of mean-elimination: ``The divisor-measure of the approximate value, added to the non-obtained quantity and the additive, being complete by two numbers with difference, then the divisor-fraction [gives] the effort. The difference of two approximate values alone [suffices]. The numerator [is] always unity [at first], thus by the intelligent. For all approximate values, the divisor-fraction is six. Thus gradually, more and more subtle approximate values are obtained.''}

\subsection*{Concluding Verses of the Bhāvitam Section (pp.\ 151--152)}

\begin{sanskritverse}
\dev{%
तथा गोले मयोक्तम् ।\\
उत्थ सदसतां मर्तानां वैराशिकमात्रमेव पाठो बुद्धिरेव बीजम् ।\\
तथा गोलाख्याय मयोक्तम् ।\\[0.3em]
अस्ति वैराशिकं पाठी बीजं च विमला मतिः ।\\
किमन्यत् सुखदीनामतो मन्दाऽर्थमुच्यते ॥\\[0.3em]
गणककर्मान्तरमयं बाललोचनगम्यं\\
सकलगणितसारं सोपपट्टिप्रकारम् ।\\
इति बहुगुणयुक्तं सर्वदोषैर्विमुक्तं\\
पठ पठ मतिबुद्ध्‌यै लाघवं प्रौढिसिद्ध्यै ॥}
\end{sanskritverse}

\cnote{The epilogue: ``Thus I have spoken of the sphere. For people ascending from the unreal to the real, the mere reading [of this text] is the seed of intelligence. Thus I have spoken of the sphere called `Gola'. There is the study of the text and pure intellect --- what else [is needed]? For the delight of the wise, this is declared. This work accessible to the eye of a child [yet] profound, the essence of all mathematics with proofs endowed, endowed with many virtues, free from all faults --- read, read with intelligence and understanding, for ease and the attainment of maturity.''}

\begin{sanskritverse}
\dev{%
इति श्रीभास्कराचार्यविरचिते सिद्धान्तशिरोमणौ\\
बीजगणिताध्यायः समाप्तः ॥}
\end{sanskritverse}

\cnote{``Thus ends the Bījagaṇita-adhyāya in the Siddhāntaśiromaṇi composed by Śrī Bhāskarācārya.''}

\bigskip

\begin{sanskritverse}
\dev{%
वि०—इति कूपाङ्कनमुखाकरे यदिह भास्करबीजमपूर्वकम् ।\\
तदुपनिषमतीव चमत्कृतिं विभिधरस्य चकार च कारणम् ॥}
\end{sanskritverse}

\cnote{Editor's note: ``That unprecedented algebra of Bhāskara which here [appears] with the face of a well-digit-mine [i.e., a numerical source], having ascertained it with the Upaniṣad [of commentary], the cause of the appearance of manifold wonder.''}

%% ========================================================================
%% Part 9: Pages 153-170
%% ========================================================================
\section{Final Sections and Publisher's Catalogue (pp.\ 153--170)}
\label{sec:final}

\subsection{Navīna-mati-viśeṣaḥ (pp.\ 153--157)}

The section on ``Special Methods for New Understanding'' presents advanced techniques for solving equations through approximation and iteration, building on the methods developed earlier.

\begin{sanskritverse}
\dev{%
नवीनमतिविशेषः ॥\\[0.5em]
निरक्ष पदं यद्‌द्विगुणात् स्यात् फलार्थं धनार्थं तदेवात्र शेषं तदस्यम् ।\\
पदद्वयं धनं शेषहृतप्रमन्यत् फलं तद्‌भवत् शेषमूलं धनेन ॥ १ ॥\\
धनार्थं नवं तस्य कृत्या विहीनो गुणः शेषमकोऽस्यवर्णस्य मानम् ।\\
मूद्‌गुसेव मते यदा शेषमानं भवेद्‌द्वितीयं तदा लव्धितौ ते ॥ २ ॥}
\end{sanskritverse}

\cnote{Rule 1: ``A number, when doubled, becomes the root; for the fruit [result], for wealth [coefficient], that same residue here --- its two terms [are] the wealth divided by the residue; the fruit [is] the root of the residue multiplied by the wealth. The new wealth is diminished by nine times its square; the multiplier, diminished [thus], [gives] the measure of the unknown. When by the opinion of Mūdgu [i.e., by the method of approximation], the residue-value becomes the second [term], then the quotient for both.''}

\subsection{Worked Examples on Special Methods (pp.\ 157--159)}

\begin{sanskritverse}
\dev{%
उदाहरणानि ॥\\[0.3em]
(१) कल्प्यते अ+क=५७६०, अ$-$क=$\frac{\text{अ}}{२}$ तदा अ=३८४०,\\
\hspace*{2em} क=१९२० इति क्षिप्रम् ॥\\[0.3em]
(२) यदि $\frac{२\text{क}+१}{२} = \frac{७\text{क}+५}{६}$ तदा कमानम्=१ इति क्षिप्रम् ॥\\[0.3em]
(३) यदि $\frac{\text{क}+१}{२} + \frac{३\text{क}-४}{५} + \frac{९}{८} = \frac{६\text{क}+७}{८}$, तदा क=२० इति क्षिप्रम् ॥\\[0.3em]
(४) $(\text{क}+\frac{५}{४})(\text{क}-\frac{५}{४}) - (\text{क}+५)(\text{क}-३) + \frac{३}{८} = ०$\\
\hspace*{2em} तदा क=१२ इति क्षिप्रम् ॥}
\end{sanskritverse}

\subsection{The Method of Approximate Roots (pp.\ 160--165)}

This section presents the \textit{āsanna-mūla} (approximate root) method --- an iterative technique for finding successively better approximations to irrational square roots, a precursor to what is now known as Newton's method.

\begin{sanskritverse}
\dev{%
आसन्नमूलानि सूत्राणि ॥\\[0.5em]
आसन्नमानस्य हराशमानम् अप्राप्तिगुण्ये सहिते क्षेपेण ।\\
पूर्णे स्याताम् अन्तरराशिकाभ्यां तदा हरांशो भवतीश्रमस्य ॥ १ ॥}
\end{sanskritverse}

The mathematical procedure is as follows: given a non-square number $n$, to find $\sqrt{n}$:
\begin{enumerate}[label=\arabic*.]
  \item Choose an initial approximation $a_1$ (the ``approximate value'')
  \item Compute the residue: $r = n - a_1^2$
  \item Form the next approximation: $a_2 = a_1 + \dfrac{r}{2a_1}$
  \item Iterate for successively better approximations
\end{enumerate}

\dev{%
कल्प्यते—\\
(१) $0$, अ, अ', अ'', $\cdots\cdots\cdots$\\
(२) $1$, सो, सो', सो'', $\cdots\cdots\cdots$\\
(३) अ, क$_१$, क$_₂$, क$_₃$, क$_₄$, $\cdots\cdots\cdots$}

\dev{%
अत्र क$_१$, क$_₂$, क$_₃$,—अयबवदासन्नमूलस्यासन्नमानानि $\frac{\text{प}}{\text{ल}}$, $\frac{\text{प}'}{\text{ल}'}$, $\frac{\text{प}''}{\text{ल}''}$ चेति ।}

$$\sqrt{n} = \frac{\frac{\sqrt{n+\text{अ}_₁}}{\text{सो}'} \cdot \text{प} + \text{प}}{\frac{\sqrt{n+\text{अ}_₁}}{\text{सो}''} \cdot \text{ल} + \text{ल}} = \frac{\text{प}(\sqrt{n+\text{अ}_₁}) + \text{सो}''\text{प}}{\text{ल}(\sqrt{n+\text{अ}_₁}) + \text{सो}''\text{ल}}$$

\subsection{Final Concluding Verses (pp.\ 166--168)}

The treatise concludes with verses that reflect on the nature of mathematical knowledge and pay homage to the tradition:

\begin{sanskritverse}
\dev{%
इति श्रीभास्कराचार्यविरचिते सिद्धान्तशिरोमणौ\\
बीजगणिताध्यायः समाप्तः ॥}
\end{sanskritverse}

\begin{sanskritverse}
\dev{%
अस्ति वैराशिकं पाठी बीजं च विमला मतिः ।\\
किमन्यत् सुखदीनामतो मन्दाऽर्थमुच्यते ॥\\[0.5em]
गणककर्मान्तरमयं बाललीलावनोगम्यं\\
सकलगणितसारं सोपपट्टिप्रकारम् ।\\
इति बहुगुणयुक्तं सर्वदोषैर्विमुक्तं\\
पठ पठ मतिबुद्ध्‌्यै लाघवं प्रौढिसिद्ध्यै ॥}
\end{sanskritverse}

\cnote{``There is the study of this work and pure intellect --- what else [is needed]? Thus is declared for the delight of the wise. This [work], accessible as a child's play [yet] the essence of all mathematics with demonstrations, endowed with many virtues, free from all faults --- read, read with intelligence and understanding, for ease and the attainment of maturity.''}

\bigskip

The author reflects on the completion of the work:

\begin{sanskritverse}
\dev{%
तथा गोले मयोक्तम् ।\\
उत्थ सदसतां मर्तानां वैराशिकमात्रमेव पाठो बुद्धिरेव बीजम् ।\\
तथा गोलाख्याय मयोक्तम् ॥}
\end{sanskritverse}

\cnote{``Thus I have spoken of the sphere. For people ascending from the unreal to the real, the mere reading [of this text] is the seed of intelligence. Thus I have spoken of what is called the sphere.''}

\subsection*{Colophon}

\begin{sanskritverse}
\dev{%
इति श्रीभास्कराचार्यविरचिते सिद्धान्तशिरोमणौ\\
बीजगणिताध्यायः समाप्तः ॥}
\end{sanskritverse}

\longline

\subsection*{Editor's Note (p.\ 168)}

\begin{sanskritverse}
\dev{%
वि०—इति कूपाङ्कनमुखाकरे यदिह भास्करबीजमपूर्वकम् ।\\
तदुपनिषमतीव चमत्कृतिं विभिधरस्य चकार च कारणम् ॥}
\end{sanskritverse}

\cnote{``That unprecedented algebra of Bhāskara, which here [appears] at the mouth of the well-digit-mine, having ascertained it with the Upaniṣad [of commentary], became the cause of the appearance of manifold wonder.''}

%% ========================================================================
%% Publisher's Catalogue (pp. 169-170)
%% ========================================================================
\section*{Publisher's Catalogue (Chowkhamba Sanskrit Book Depot)}
\addcontentsline{toc}{section}{Publisher's Catalogue}

The final two pages contain a catalogue of Sanskrit publications available from the \textsc{Chowkhamba Sanskrit Book Depot} (Chowkhamba Griham), 4015 Thatheri Bazar, Varanasi. The catalogue lists works across all branches of Sanskrit learning, including:

\begin{longtable}{r p{8cm} c c}
\toprule
\textbf{No.} & \textbf{Title} & \textbf{Price} & \textbf{Pages} \\
\midrule
\endhead
19 & Golaprakāśasya Golādhyāya-Bījagaṇita-vāsiṣṭha-\\
   & koṇa-gaṇite Śrī-Candraśekhāra-saṃskṛtite & $\cdots$ & ८ & ७ \\
20 & Golaprakāśīya-cāpīyatrikōṇa-gaṇitasya\\
   & ``Vivāha'' ṭīkā Śrīmahāvīra-paṇḍita-viracitā & $\cdots$ & ७ & ८ \\
21 & Vāsanāvyañjana-ṭīkā Śrīgaṇeśa-dhara-mitra-kṛtā\\
   & ṭīkā-sahitā & $\cdots$ & ९ & ७ \\
22 & Prasaṅga-mūla-prabhāvicāra-syāt-mūla-prathः & $\cdots$ & ७ & ४५ \\
23 & Karṇa-prakāśaḥ Śrī-sudhā-kara-dvivedi-saṃśodhitः & $\cdots$ & १ & ८ \\
24 & Jātaka-paddhatiḥ & $\cdots$ & २ & ७ \\
25 & Bhāvate-karaṇaṃ (Bījakoṣa-udāharaṇa-sahitam) & $\cdots$ & २ & ७ \\
26 & Prabhāve-padavः & $\cdots$ & $\cdots$ & ७ \\
27 & Daivajña-kāma-dhenuः & $\cdots$ & ४ & ८ \\
28 & Sūrya-siddhāntaṃ (Mṛtaṃ) Śrī-sudhā-kara-dvivedi-\\
   & kṛta-ṭīkopetaḥ & $\cdots$ & ३ & ७ \\
29 & Makaranda-vivaraṇam (udāharaṇa-sahitam) & $\cdots$ & ७ & ६ \\
30 & Jātaka-paddhatiḥ proktamārga-vyākhyā-sahitā & $\cdots$ & ७ & १२ \\
31 & Digvijaya-saṃ M.M. Śrī-sudhā-kara-dvivedi-viracitā & $\cdots$ & ७ & १० \\
32 & Gaṇaka-taraṅgiṇī Śrī-sudhā-kara-dvivedi-nirmitā & $\cdots$ & ३ & ७ \\
33 & Yājūṣa-jyotiṣam (Śrīmākara-sudhā-kara-bhāṣya-sahitam) & $\cdots$ & १ & ४ \\
34 & Piṇḍa-prabhākaraḥ Śrī-sudhā-kara-dvivedi-viracitaḥ & $\cdots$ & ७ & ४ \\
35 & Suktiratnāvaliḥ (Mṛtaṃ) Śrī-sudhā-kara-dvivedi-viracitः & $\cdots$ & ७ & ७ \\
36 & Vallana-kalanam—Śrīmuralī-dhara-dvivedi-viracitam\\
   & (caturtha-adhyāya-paryantam) & $\cdots$ & १ & ७ \\
37 & Gaṇitakā itihāsa Pāṭīgaṇita Śrī-sudhā-kara-\\
   & dvivedi viracitam & $\cdots$ & २ & ७ \\
38 & Bṛhat-saṃhitā (Mahārāja-ṭīkā-sahitā) & $\cdots$ & १६ & ४ \\
39 & Nāradī-saṃhitā & $\cdots$ & ७ & ८ \\
40 & Vāsatvavibodha-prabhāsa-saṃgrahaḥ (Mṛtaṃ)\\
   & Śrī-sudhā-kara-dvivedi-viracitaḥ & $\cdots$ & ७ & २ \\
41 & Tulanātmaka bhāṣā śāstra—arthavāda—bhāṣā viṣayan\\
   & (lekhaṭa—ḍākṭara Maṅgala-deva śāstrī)\ldots & २ & १० \\
\bottomrule
\end{longtable}

\vfill
\begin{center}
\textbf{\dev{कृष्णदास गुप्त,}}\\
\textbf{४०१५ ठठेरी बाजार, बनारस सिटी ।}\\[0.5em]
\textsc{Chowkhamba Sanskrit Book Depot}\\
4015 Thatheri Bazar, Varanasi City
\end{center}

%% ========================================================================
%% End of Document
%% ========================================================================
\end{document}
